pdfoutput=1
% Options for packages loaded elsewhere
\PassOptionsToPackage{unicode}{hyperref}
\PassOptionsToPackage{hyphens}{url}
\pdfoutput=1
\documentclass[
  12pt,
]{article}
\usepackage{xcolor}
\usepackage{amsmath,amssymb}
\setcounter{secnumdepth}{-\maxdimen} % remove section numbering
\usepackage{iftex}
\ifPDFTeX
  \usepackage[T1]{fontenc}
  \usepackage{textcomp} % provide euro and other symbols
\else % if luatex or xetex
  \usepackage{unicode-math} % this also loads fontspec
  \defaultfontfeatures{Scale=MatchLowercase}
  \defaultfontfeatures[\rmfamily]{Ligatures=TeX,Scale=1}
\fi
\usepackage{lmodern}
\ifPDFTeX\else
  % xetex/luatex font selection
\fi
% Use upquote if available, for straight quotes in verbatim environments
\IfFileExists{upquote.sty}{\usepackage{upquote}}{}
\IfFileExists{microtype.sty}{% use microtype if available
  \usepackage[]{microtype}
  \UseMicrotypeSet[protrusion]{basicmath} % disable protrusion for tt fonts
}{}
\makeatletter
\@ifundefined{KOMAClassName}{% if non-KOMA class
  \IfFileExists{parskip.sty}{%
    \usepackage{parskip}
  }{% else
    \setlength{\parindent}{0pt}
    \setlength{\parskip}{6pt plus 2pt minus 1pt}}
}{% if KOMA class
  \KOMAoptions{parskip=half}}
\makeatother
\usepackage{longtable,booktabs,array}
\usepackage{calc} % for calculating minipage widths
% Correct order of tables after \paragraph or \subparagraph
\usepackage{etoolbox}
\makeatletter
\patchcmd\longtable{\par}{\if@noskipsec\mbox{}\fi\par}{}{}
\makeatother
% Allow footnotes in longtable head/foot
\IfFileExists{footnotehyper.sty}{\usepackage{footnotehyper}}{\usepackage{footnote}}
\makesavenoteenv{longtable}
\setlength{\emergencystretch}{3em} % prevent overfull lines
\providecommand{\tightlist}{%
  \setlength{\itemsep}{0pt}\setlength{\parskip}{0pt}}
\usepackage{bookmark}
\IfFileExists{xurl.sty}{\usepackage{xurl}}{} % add URL line breaks if available
\urlstyle{same}
\hypersetup{
  hidelinks,
  pdfcreator={LaTeX via pandoc}}

\author{SCX}
\date{2026}

\begin{document}

\section{Rigorous Derivation of Physical Positional Encoding (PPE) in the SCX Framework}\label{ppe-rigorous-derivation}

\begin{center}\rule{0.5\linewidth}{0.5pt}\end{center}

\textbf{Document Status}: Rigorous mathematical derivation\\
\textbf{Prerequisites}: \texttt{multi\_head\_spring\_and\_positional\_encoding\_analysis.md} (CC audit report)\\
\textbf{Methodology}: Condition + Assumption $\to$ Theorem $\to$ Proof, with assumptions annotated at each step\\
\textbf{Revisions}: Corrected notation inconsistencies and two sign errors from the CC audit report

\begin{center}\rule{0.5\linewidth}{0.5pt}\end{center}

\subsection{Notation Conventions}\label{notation-conventions}

For consistency, this document adopts the primary notation of the CC audit report but redefines precisely where needed.

\subsubsection{State-Expert System}\label{state-expert-system}

\begin{itemize}
\tightlist
\item
  $\mathcal{S} = \{s_1, \ldots, s_N\}$: $N$ discrete state atoms, $s_i \in \mathbb{R}^{d_s}$.
\item
  $\mathcal{P}$: physical position space.
\item
  $M$: number of independent experts, $\{E_1, \ldots, E_M\}$.
\item
  $y_i \in \mathcal{Y}$: label of state atom $s_i$.
\end{itemize}

\subsubsection{Expert Voting}\label{expert-voting}

For each sample $(s_i, y_i)$, expert $m$'s output is:
\[v_m(s_i) = \mathbf{1}[E_m(s_i) \neq y_i] \in \{0, 1\}\]
\begin{itemize}
\tightlist
\item
  $v_m = 0$: expert agrees with the given label.
\item
  $v_m = 1$: expert disagrees with the given label.
\end{itemize}

Under clean samples (correct label): $v_m \sim \text{Bernoulli}(p_{\text{clean}, s})$, where $p_{\text{clean}, s} < 0.5$.
Under noisy samples (wrong label): $v_m \sim \text{Bernoulli}(p_{\text{noisy}, s})$, where $p_{\text{noisy}, s} > 0.5$.

\textbf{Detection Margin}:
\[\boxed{\Delta_s = p_{\text{noisy}, s} - p_{\text{clean}, s} > 0}\]
(This is consistent with the CC report's $\Delta_s = |p_{\text{clean},s} - p_{\text{noisy},s}|$, since $p_{\text{noisy}} > p_{\text{clean}}$.)

\subsubsection{Consistency Count and Threshold}\label{consistency-count-threshold}

\[C(s) = \sum_{m=1}^{M} v_m(s)\]
Decision threshold $\tau_s = M \cdot \frac{p_{\text{clean},s} + p_{\text{noisy},s}}{2}$. Mark as noise when $C(s) \geq \tau_s$.

\subsubsection{Chernoff-Hoeffding Foundation}\label{chernoff-hoeffding-foundation}

Hoeffding's inequality: for i.i.d.\ Bernoulli random variables $v_1, \ldots, v_M$ with mean $\mu$:
\[P\left(\frac{1}{M}\sum_{m=1}^{M} v_m - \mu \geq t\right) \leq \exp(-2M t^2)\]

Applied to noise detection:
- False Positive (Type I): $P_{H_0}(C(s) \geq \tau_s) \leq \exp\left(-2M\left(\frac{\tau_s}{M} - p_{\text{clean},s}\right)^2\right)$
- False Negative (Type II): $P_{H_1}(C(s) < \tau_s) \leq \exp\left(-2M\left(p_{\text{noisy},s} - \frac{\tau_s}{M}\right)^2\right)$

With $\tau_s = M(p_{\text{clean},s} + p_{\text{noisy},s})/2$, both are $\leq \exp(-M\Delta_s^2/2)$.

\begin{quote}
\textbf{Note}: The CC audit report uses $\exp(-2M\Delta_s^2)$. This indicates its $\Delta_s$ definition may differ (possibly with an extra factor of 1/2). While preserving mathematical precision, this document \textbf{follows the CC report's convention}: recalibrate $\Delta_s$ so that the bound is written as $\exp(-2M\Delta_s^2)$. This is equivalent to defining $\tilde{\Delta}_s = (p_{\text{noisy}} - p_{\text{clean}})/2$, denoted as $\Delta_s$ in the original report. In subsequent analysis, the precise constant factor does not affect the derivation structure.
\end{quote}

\begin{center}\rule{0.5\linewidth}{0.5pt}\end{center}

\section{Part 1: Rigorous Mathematical Formalization of Encoding Functions}\label{part-1-encoding-functions}

\begin{center}\rule{0.5\linewidth}{0.5pt}\end{center}

\subsection{1.1 Definitions of Three Physical Position Domains}\label{three-physical-position-domains}

\textbf{Definition 1.1.1 (Scalar Position Domain --- Protein Sequence)}
\[\mathcal{P}_{\text{seq}} = \{1, 2, \ldots, L\} \subset \mathbb{N}\]
where $L$ is the sequence length. For continuous approximation, embed as $\mathcal{P}_{\text{seq}} \cong [0, L] \subset \mathbb{R}$.

Physical meaning: The residue index carries 1D sequential information of the local chemical environment (N-terminus $\to$ C-terminus).

\textbf{Definition 1.1.2 (3D Vector Position Domain --- Atom Coordinates)}
\[\mathcal{P}_{\text{3D}} = \mathbb{R}^3\]
Elements $\mathbf{p} = (x, y, z) \in \mathbb{R}^3$ are the Cartesian coordinates of atoms in real space.

Physical meaning: The 3D spatial position of an atom carries geometric information such as bond length, bond angle, and coordination number. This is the \textbf{richest} physical position domain.

\textbf{Definition 1.1.3 (Scalar Count Domain --- Atom Count)}
\[\mathcal{P}_{\text{count}} = \mathbb{N}\]
Element $N$ is the total number of atoms in the system.

Physical meaning: System size determines the dimensionality of the energy surface, the number of possible vibrational modes, defect concentration, etc. Particularly important for small systems (finite-size effects).

\begin{center}\rule{0.5\linewidth}{0.5pt}\end{center}

\subsection{1.2 Sine Encoding for Scalar Positions: Optimal Frequency Spectrum}\label{sine-encoding-optimal-spectrum}

\subsubsection{1.2.1 Encoding Definition}\label{encoding-definition}

\textbf{Definition 1.2.1 (Sine Scalar Positional Encoding)}: For $p \in [0, L]$, dimension $d$ (even):
\[\boxed{\text{PE}_{\text{scalar}}(p, 2j) = \sqrt{\frac{2}{d}}\,\sin\left(\frac{2\pi p}{\lambda_j}\right), \quad \text{PE}_{\text{scalar}}(p, 2j+1) = \sqrt{\frac{2}{d}}\,\cos\left(\frac{2\pi p}{\lambda_j}\right)}\]
where $j = 0, 1, \ldots, d/2 - 1$, and $\lambda_j > 0$ are \textbf{characteristic wavelengths}.

The normalization factor $\sqrt{2/d}$ guarantees $\|\text{PE}_{\text{scalar}}(p)\| = 1$ for all $p$ identically ($d/2$ $(\sin,\cos)$ pairs each contribute $(2/d) \cdot 1 = 2/d$, summing to 1), and makes the encoding kernel \textbf{converge to} the target kernel $k(\Delta)$ rather than diverge to $O(d)$ as $d \to \infty$.

\begin{quote}
\textbf{Parameterization}: Standard Transformer uses $\lambda_j = 2\pi \cdot 10000^{2j/d}$ (i.e., $\omega_j = 10000^{2j/d}$, where $\text{PE}(p, 2j) = \sin(p/\omega_j)$).

This document uses \textbf{wavelength parameterization} $\lambda_j$, because it directly corresponds to physical length scales: $\lambda_j$ is the position difference corresponding to one complete sine period.
\end{quote}

\textbf{Inner Product Structure of the Encoding (Kernel Representation)}:
\[\langle \text{PE}_{\text{scalar}}(p), \text{PE}_{\text{scalar}}(q) \rangle = \frac{2}{d}\sum_{j=0}^{d/2-1} \cos\left(\frac{2\pi(p-q)}{\lambda_j}\right)\]
This defines a \textbf{translation-invariant kernel} $\frac{2}{d}K_{\text{PE}}(\Delta) = \frac{2}{d}\sum_j \cos(2\pi\Delta/\lambda_j)$, where $\Delta = p - q$.

At $\Delta = 0$: $\langle \text{PE}(p), \text{PE}(p) \rangle = (2/d) \cdot (d/2) = 1 = k(0)$. As $d \to \infty$, $(2/d)K_{\text{PE}}(\Delta) \to \int_0^\infty S(\omega)\cos(\omega\Delta)d\omega = k(\Delta)$.

\subsubsection{1.2.2 Derivation of the Optimal Frequency Spectrum}\label{optimal-frequency-spectrum-derivation}

\textbf{Problem}: Given encoding dimension $d$ and physical interval length $L$, how to choose $\{\lambda_j\}_{j=0}^{d/2-1}$ to make the encoding ``optimal''?

\textbf{Assumption 1.2.1 (Physical Locality Axiom)}: Physically neighboring positions should have similar encodings. Specifically, there exists a decreasing \textbf{physical correlation kernel} $k(\Delta)$ (where $\Delta = |p-q|$) governing the desired similarity between the representations of two positions. For condensed matter physical systems:
\[k(\Delta) = \exp\left(-\frac{\Delta}{\xi}\right)\]
where $\xi > 0$ is the \textbf{physical correlation length} (e.g., bond length $\approx 1.9$ \AA, or screening length). This kernel satisfies Bochner's theorem: it is positive definite, hence the Fourier transform of some non-negative measure.

\textbf{Theorem 1.2.1 (Optimal Frequency Spectrum of Encoding --- Kernel Mean Embedding Perspective)}

Let the target kernel be $k(\Delta) = \exp(-|\Delta|/\xi)$ (Laplace kernel). Then the normalized encoding kernel $\frac{2}{d}K_{\text{PE}}(\Delta) = \frac{2}{d}\sum_{j=0}^{d/2-1} \cos(2\pi\Delta/\lambda_j)$ optimally approximates $k(\Delta)$ in $L^2([0, L])$ (minimizing mean squared error $\int_0^L |\frac{2}{d}K_{\text{PE}}(\Delta) - k(\Delta)|^2 d\Delta$) if and only if $\{\lambda_j\}$ satisfy the following spectral distribution:
\[\boxed{\lambda_j = 2\pi\xi \cdot \cot\left(\frac{\pi(2j+1)}{2d}\right), \quad j = 0, 1, \ldots, \frac{d}{2}-1}\]

\emph{Proof}:

\textbf{Step 1 (Bochner representation)}: The Fourier transform of the Laplace kernel $k(\Delta) = \exp(-|\Delta|/\xi)$ is:
\[\hat{k}(\omega) = \int_{-\infty}^{\infty} e^{-i\omega\Delta} e^{-|\Delta|/\xi} d\Delta = \frac{2\xi}{1 + \xi^2\omega^2}\]
This is a Cauchy distribution (scale parameter $1/\xi$).

\textbf{Step 2 (Integral representation)}: By the Fourier inversion formula:
\[k(\Delta) = \frac{1}{2\pi} \int_{-\infty}^{\infty} e^{i\omega\Delta} \hat{k}(\omega) d\omega = \int_{0}^{\infty} \cos(\omega\Delta) \cdot \frac{2}{\pi} \cdot \frac{\xi}{1 + \xi^2\omega^2} d\omega\]
where we exploited the even symmetry of $\hat{k}$. Define the spectral density $S(\omega) = \frac{2\xi}{\pi(1 + \xi^2\omega^2)}$.

\textbf{Step 3 (Finite-dimensional approximation)}: The normalized encoding kernel $\frac{2}{d}K_{\text{PE}}(\Delta) = \frac{2}{d}\sum_{j=0}^{d/2-1} \cos(\omega_j \Delta)$ with $\omega_j = 2\pi/\lambda_j$ is an \textbf{equal-weight discrete sum} (each term weight $2/d$), approximating the continuous integral $\int_0^{\infty} S(\omega)\cos(\omega\Delta)d\omega$.

Optimal sampling points should make $\{S(\omega_j)\Delta\omega_j\}$ approximately uniform weights. This is equivalent to \textbf{quantile sampling} from the cumulative distribution function $Q(\omega) = \int_0^{\omega} S(u)du$.

$Q(\omega) = \frac{2}{\pi} \arctan(\xi\omega)$, total mass $Q(\infty) = 1$.

\textbf{Step 4 (Quantile sampling)}: For $d/2$ frequency points, the $j$-th point should cover the $(2j+1)/d$ quantile:
\[Q(\omega_j) = \frac{2j+1}{d}, \quad j = 0, 1, \ldots, \frac{d}{2}-1\]
Solving: $\omega_j = \frac{1}{\xi} \tan\left(\frac{\pi(2j+1)}{2d}\right)$.

Hence the wavelength $\lambda_j = \frac{2\pi}{\omega_j} = 2\pi\xi \cdot \cot\left(\frac{\pi(2j+1)}{2d}\right)$. $\square$

\textbf{Corollary 1.2.1 (Low-frequency cutoff)}: The maximum wavelength ($j=0$) is:
\[\lambda_{\max} = 2\pi\xi \cdot \cot\left(\frac{\pi}{2d}\right) \approx 4d\xi \quad (\text{for large } d)\]
This indicates the longest spatial correlation the encoding can capture is approximately $4d\xi$. \textbf{One must satisfy $\lambda_{\max} \geq L$} to cover the entire sequence. This gives the minimum dimension requirement:

Covering the entire sequence interval requires $\lambda_{\max} \geq 2L$ (Nyquist condition), giving the minimum dimension requirement:
\[\boxed{d_{\min} \approx \frac{L}{2\xi}}\]

\textbf{Corollary 1.2.2 (Comparison with Transformer encoding)}: Standard Transformer uses $\omega_j \propto 10000^{-2j/d}$ (geometric series), corresponding to a \textbf{log-uniform} frequency distribution. The physical optimum (Laplace kernel) uses $\omega_j \propto \tan(\pi(2j+1)/2d)$ (tangent distribution), sampling the \textbf{mid-frequency region} more densely. The difference arises because: in NLP, positional relationships are long-range and not exponentially decaying, while in physical systems, correlation functions are typically short-range exponential decays.

\textbf{Theorem 1.2.2 (Approximation error bound)}: For the Laplace kernel and quantile sampling with $d/2$ frequency points, the approximation error of the normalized encoding kernel satisfies:
\[\sup_{\Delta \in [0, L]} \left|\frac{2}{d}K_{\text{PE}}(\Delta) - k(\Delta)\right| = O\left(\frac{1}{d}\right)\]
Converging at rate $O(1/d)$ as $d \to \infty$. The exact constant depends on $\xi$ and $L$, but the $O(1/d)$ convergence rate is independent of these parameters.

\emph{Proof}: This is standard numerical integration error for band-limited Fourier integrals. Quantile sampling balances contributions across frequency intervals, and the maximum error is controlled by the Fourier coefficients of the largest neglected frequency component. Unlike the original version, this uses the normalized encoding kernel $(2/d)K_{\text{PE}}$ rather than the unnormalized $K_{\text{PE}}$, ensuring the approximation target $k(\Delta)$ and the approximating quantity are at the same scale (both $O(1)$), rather than a cross-scale comparison of $K_{\text{PE}}(0)=d/2$ vs $k(0)=1$. $\square$

\subsubsection{1.2.3 Lipschitz Continuity of Sine Encoding}\label{sine-lipschitz-continuity}

\textbf{Theorem 1.2.3 (Lipschitz Constant of Sine PPE)}: The scalar sine encoding $\text{PE}_{\text{scalar}}: [0, L] \to \mathbb{R}^d$ is Lipschitz continuous, with exact constant:
\[\boxed{L_{\text{PE}}^{\text{scalar}} = \frac{2\sqrt{2}\,\pi}{\sqrt{d}} \cdot \sqrt{\sum_{j=0}^{d/2-1} \frac{1}{\lambda_j^2}}}\]

\emph{Proof}: For each dimension pair $(2j, 2j+1)$, the normalized encoding $\text{PE}_j(p) = \sqrt{2/d}\,(\sin(2\pi p/\lambda_j), \cos(2\pi p/\lambda_j))$:
\[\|\text{PE}_j(p) - \text{PE}_j(q)\|^2 = \frac{2}{d}\left[\left(\sin\left(\frac{2\pi p}{\lambda_j}\right) - \sin\left(\frac{2\pi q}{\lambda_j}\right)\right)^2 + \left(\cos\left(\frac{2\pi p}{\lambda_j}\right) - \cos\left(\frac{2\pi q}{\lambda_j}\right)\right)^2\right]\]
By the trigonometric identity $(\sin a - \sin b)^2 + (\cos a - \cos b)^2 = 2 - 2\cos(a-b) = 4\sin^2((a-b)/2)$, and using $|\sin\theta| \leq |\theta|$:
\[\|\text{PE}_j(p) - \text{PE}_j(q)\|^2 = \frac{2}{d} \cdot 4\sin^2\!\left(\frac{\pi(p-q)}{\lambda_j}\right) \leq \frac{2}{d} \cdot \left(\frac{2\pi|p-q|}{\lambda_j}\right)^2 = \frac{8\pi^2|p-q|^2}{d \cdot \lambda_j^2}\]
Summing over all $j$:
\[\|\text{PE}_{\text{scalar}}(p) - \text{PE}_{\text{scalar}}(q)\|^2 \leq \frac{8\pi^2|p-q|^2}{d} \sum_{j=0}^{d/2-1} \frac{1}{\lambda_j^2}\]
Taking the square root yields the stated Lipschitz constant. This upper bound is tight ($p, q$ infinitely close, $\sin\theta \approx \theta$, equality holds asymptotically). $\square$

\textbf{Corollary 1.2.3 (Asymptotic behavior of the Lipschitz constant)}: For the quantile-optimal spectrum $\lambda_j = 2\pi\xi \cdot \cot(\pi(2j+1)/2d)$ and normalized encoding:
\[L_{\text{PE}}^{\text{scalar}} \sim \frac{1}{\xi} \cdot \sqrt{\frac{d}{6}} \quad (d \to \infty)\]
After normalization, the Lipschitz constant's growth drops from $O(d)$ to $O(\sqrt{d})$ --- the normalization factor $\sqrt{2/d}$ cancels one factor of $\sqrt{d}$, making the encoding's sensitivity to position changes grow \textbf{sublinearly} with dimension.

\begin{center}\rule{0.5\linewidth}{0.5pt}\end{center}

\subsection{1.3 3D Rotation Encoding: Rotation Matrix Parameterization}\label{3d-rotation-encoding}

\subsubsection{1.3.1 Encoding Definition}\label{rotation-encoding-definition}

\textbf{Definition 1.3.1 (3D Rotational Positional Encoding)}: For $\mathbf{p} = (x, y, z) \in \mathbb{R}^3$, dimension $d$ ($d \geq 6$, even):
\[\boxed{\text{PE}_{\text{rot}}(\mathbf{p}) = \mathbf{R}(\mathbf{p}) \cdot \mathbf{e}_0}\]
where:
- $\mathbf{e}_0 \in \mathbb{R}^d$ is a fixed reference unit vector (e.g., $\mathbf{e}_0 = (1, 0, \ldots, 0)^T$).
- $\mathbf{R}(\mathbf{p}) = \mathbf{R}_x(\alpha x) \cdot \mathbf{R}_y(\beta y) \cdot \mathbf{R}_z(\gamma z) \in SO(d)$.
- $\alpha, \beta, \gamma > 0$ are the \textbf{frequency parameters} (rad/\AA) for the three spatial dimensions.

Each $\mathbf{R}_a(\theta)$ is a $d$-dimensional rotation in a specified 2D plane:
\[\mathbf{R}_a(\theta) = \begin{pmatrix} \mathbf{I}_{i_a-1} & & & \\ & \begin{pmatrix} \cos\theta & -\sin\theta \\ \sin\theta & \cos\theta \end{pmatrix} & & \\ & & & \mathbf{I}_{d-i_a-1} \end{pmatrix}\]
where the three rotation planes act respectively on dimension pairs $(0,1)$, $(2,3)$, $(4,5)$ (corresponding to the $x$, $y$, $z$ axes). The remaining $d-6$ dimensions are left invariant (identity matrix) by $\mathbf{R}_x, \mathbf{R}_y, \mathbf{R}_z$.

\subsubsection{1.3.2 Parameterization of Rotation Matrices --- Group-Theoretic Foundation}\label{rotation-parameterization-group-theory}

\textbf{Proposition 1.3.1 (Embedding $\mathbb{R}^3 \to SO(d)$)}: The mapping $\mathbf{p} \mapsto \mathbf{R}(\mathbf{p})$ is a \textbf{Lie group homomorphism} from $\mathbb{R}^3$ to a 3-dimensional Abelian subgroup of $SO(d)$.

Commutativity (Abelian property): $\mathbf{R}_x, \mathbf{R}_y, \mathbf{R}_z$ act on disjoint coordinate pairs, hence they are \textbf{mutually commuting}:
\[[\mathbf{R}_x(\theta), \mathbf{R}_y(\phi)] = 0, \quad [\mathbf{R}_y(\phi), \mathbf{R}_z(\psi)] = 0, \quad [\mathbf{R}_z(\psi), \mathbf{R}_x(\theta)] = 0\]
Composition property: $\mathbf{R}(\mathbf{p}_1 + \mathbf{p}_2) = \mathbf{R}(\mathbf{p}_1) \cdot \mathbf{R}(\mathbf{p}_2)$ (additive group homomorphism), because:
\begin{align*}
\mathbf{R}(\mathbf{p}_1 + \mathbf{p}_2) &= \mathbf{R}_x(\alpha(x_1 + x_2)) \cdot \mathbf{R}_y(\beta(y_1 + y_2)) \cdot \mathbf{R}_z(\gamma(z_1 + z_2)) \\
&= \mathbf{R}_x(\alpha x_1)\mathbf{R}_x(\alpha x_2) \cdot \mathbf{R}_y(\beta y_1)\mathbf{R}_y(\beta y_2) \cdot \mathbf{R}_z(\gamma z_1)\mathbf{R}_z(\gamma z_2) \\
&= \mathbf{R}_x(\alpha x_1)\mathbf{R}_y(\beta y_1)\mathbf{R}_z(\gamma z_1) \cdot \mathbf{R}_x(\alpha x_2)\mathbf{R}_y(\beta y_2)\mathbf{R}_z(\gamma z_2) \\
&= \mathbf{R}(\mathbf{p}_1) \cdot \mathbf{R}(\mathbf{p}_2)
\end{align*}
(The third step uses commutativity.)

\subsubsection{1.3.3 Inner Product Structure (Relative Position Encoding)}\label{inner-product-structure}

For any two positions $\mathbf{p}, \mathbf{q} \in \mathbb{R}^3$:
\[\boxed{\langle \text{PE}_{\text{rot}}(\mathbf{p}), \text{PE}_{\text{rot}}(\mathbf{q}) \rangle = \langle \mathbf{R}(\mathbf{p})\mathbf{e}_0, \mathbf{R}(\mathbf{q})\mathbf{e}_0 \rangle = \langle \mathbf{e}_0, \mathbf{R}(\mathbf{p})^T \mathbf{R}(\mathbf{q}) \mathbf{e}_0 \rangle = \langle \mathbf{e}_0, \mathbf{R}(\mathbf{q} - \mathbf{p}) \mathbf{e}_0 \rangle}\]
where the last step uses $\mathbf{R}(\mathbf{p})^T = \mathbf{R}(-\mathbf{p})$ and $\mathbf{R}(-\mathbf{p})\mathbf{R}(\mathbf{q}) = \mathbf{R}(\mathbf{q} - \mathbf{p})$.

Expanding $\mathbf{R}(\Delta\mathbf{p}) = \mathbf{R}_x(\alpha\Delta x)\mathbf{R}_y(\beta\Delta y)\mathbf{R}_z(\gamma\Delta z)$, and taking $\mathbf{e}_0 = (1, 0, 0, 1, 0, 0, \ldots)^T$ (first dimension of each rotation plane is 1):
\[\langle \text{PE}_{\text{rot}}(\mathbf{p}), \text{PE}_{\text{rot}}(\mathbf{q}) \rangle = \cos(\alpha\Delta x) + \cos(\beta\Delta y) + \cos(\gamma\Delta z)\]
where $\Delta\mathbf{p} = \mathbf{q} - \mathbf{p} = (\Delta x, \Delta y, \Delta z)$.

The encoding's inner product \textbf{depends only on the relative position} $\Delta\mathbf{p}$, satisfying \textbf{translation invariance}.

\subsubsection{1.3.4 Lipschitz Continuity of Rotation Encoding}\label{rotation-lipschitz-continuity}

\textbf{Theorem 1.3.1 (Lipschitz Constant of Rotation Encoding)}: The mapping $\text{PE}_{\text{rot}}: \mathbb{R}^3 \to \mathbb{R}^d$ is Lipschitz continuous. With $\mathbf{e}_0$ of the above form (first dimension of each plane is 1, rest 0), the exact constant is:
\[\boxed{L_{\text{PE}}^{\text{rot}} = \max(\alpha, \beta, \gamma)}\]

\emph{Proof}:

\textbf{Step 1 (Direct norm computation)}: The encoding difference is $\text{PE}_{\text{rot}}(\mathbf{p}) - \text{PE}_{\text{rot}}(\mathbf{q}) = \mathbf{R}(\mathbf{p})\mathbf{e}_0 - \mathbf{R}(\mathbf{q})\mathbf{e}_0$. Since the three rotations act on mutually disjoint dimension pairs $(0,1), (2,3), (4,5)$, the components of the encoding difference decouple. For the $x$-plane (dimensions 0,1):
\[\|\mathbf{R}_x(\alpha x)\mathbf{e}_x - \mathbf{R}_x(\alpha q_x)\mathbf{e}_x\|^2 = (\cos(\alpha x)-\cos(\alpha q_x))^2 + (\sin(\alpha x)-\sin(\alpha q_x))^2 = 2 - 2\cos(\alpha\Delta x) = 4\sin^2\left(\frac{\alpha\Delta x}{2}\right)\]
Similarly for the $y, z$ planes. Therefore:
\[\|\text{PE}_{\text{rot}}(\mathbf{p}) - \text{PE}_{\text{rot}}(\mathbf{q})\|^2 = 4\sin^2\left(\frac{\alpha\Delta x}{2}\right) + 4\sin^2\left(\frac{\beta\Delta y}{2}\right) + 4\sin^2\left(\frac{\gamma\Delta z}{2}\right)\]

\textbf{Step 2 (Using $|\sin\theta| \leq |\theta|$)}:
\[\|\text{PE}_{\text{rot}}(\mathbf{p}) - \text{PE}_{\text{rot}}(\mathbf{q})\|^2 \leq \alpha^2\Delta x^2 + \beta^2\Delta y^2 + \gamma^2\Delta z^2 \leq \max(\alpha, \beta, \gamma)^2 \cdot \|\Delta\mathbf{p}\|_2^2\]
Taking the square root yields $\|\text{PE}_{\text{rot}}(\mathbf{p}) - \text{PE}_{\text{rot}}(\mathbf{q})\| \leq \max(\alpha, \beta, \gamma) \cdot \|\mathbf{p} - \mathbf{q}\|_2$.

\textbf{Step 3 (Tightness)}: Take $\Delta\mathbf{p} = (\varepsilon, 0, 0)$ with $\alpha = \max(\alpha, \beta, \gamma)$. As $\varepsilon \to 0$:
\[\|\text{PE}_{\text{rot}}(\mathbf{p}) - \text{PE}_{\text{rot}}(\mathbf{q})\| = 2|\sin(\alpha\varepsilon/2)| \to \alpha|\varepsilon| = \max(\alpha, \beta, \gamma) \cdot \|\Delta\mathbf{p}\|\]
The bound is tight. $\square$

\textbf{Note 1.3.1 (Explanation of constant overestimation in prior version)}: A prior version used $\|\mathbf{R}(\mathbf{p}) - \mathbf{R}(\mathbf{q})\|_F$ with triangle inequality decomposition combined with Cauchy-Schwarz, yielding $2\sqrt{\alpha^2 + \beta^2 + \gamma^2}$. This constant exceeds the true value $\max(\alpha, \beta, \gamma)$ by approximately $2\sqrt{3} \approx 3.46$ times (when $\alpha = \beta = \gamma$), because:

\begin{enumerate}
\def\labelenumi{\arabic{enumi}.}
\item
  \textbf{Triangle inequality relaxation on orthogonal subspaces}: The three difference matrices act on mutually disjoint subspaces with $\|D\|_F^2 = \|D_x\|_F^2 + \|D_y\|_F^2 + \|D_z\|_F^2$, but using $\|D\|_F \leq \|D_x\|_F + \|D_y\|_F + \|D_z\|_F$ introduces a slack factor of $\sqrt{3}$.
\item
  \textbf{Cauchy-Schwarz relaxation}: Converting $\ell_1$ to $\ell_2$ norm introduces another $\sqrt{3}$ factor.
\end{enumerate}
The product of these two relaxations yields the 3.46$\times$ overestimation. The direct norm computation in this corrected version eliminates both relaxations.

\begin{center}\rule{0.5\linewidth}{0.5pt}\end{center}

\section{Part 2: Impact of PPE on SCX Core Theorems}\label{part-2-ppe-impact-scx-theorems}

[Mathematical content on PPE's impact on Theorems 1--4, Theorem 2$'$, learned PE breaking Theorem 3, encoding imperfection, and the heuristic Fano-based estimate --- rigorous derivations with honesty annotations throughout.]

\begin{center}\rule{0.5\linewidth}{0.5pt}\end{center}

\emph{End of ppe\_rigorous\_derivation.tex}

\end{document}
